\chapter{Исследовательская часть}
%<цель исследования, на какой машине делали (указать ЦПУ, ОЗУ, ОС), желательно написать, как исследовали и при каких условиях; что получили в результате (таблицы + графики)>

\section{Замеры процессорного времени}

Замеры проводились на ноутбуке Redmibook Pro 14 2022:
\begin{itemize}
	\item процессор Intel i7-12650H (архитектура x86-64, 10 ядер по 2.3 Ггц);
	\item оперативная память 16 Гб (DDR4, 5200 МГц);
	\item операционная система Windows 11 Home 25H2.
\end{itemize}
Для замеров использовалась функция clock() из библиотеки ctime~\cite[с.~288]{lit2}.

% выравнивание по точке
\begin{table}[!htbp]
	\caption{Некоторые замеры процессорного времени для различных N (при k = 1000).}
	\centering
	\begin{tabularx}{\textwidth}{|>{\centering\arraybackslash}X|*{3}{>{\centering\arraybackslash}S[table-format=4.5]|}}
		\hline
		\textbf{N} & \textbf{\hspace{3em}Рекурсивный\hspace{3em}} & \textbf{\hspace{3em}Нерекурсивный\hspace{3em}} \\\hline
		1 & 0.067 & 0.051 \\\hline
		10 & 0.678 & 0.717 \\\hline
		20 & 1.523 & 1.912 \\\hline
		30 & 2.081 & 2.1885 \\\hline
	\end{tabularx}
\end{table}

% график должен быть читаем в чб
% графики иллюстрируют

% таблица + график из txt/csv

\begin{figure}[!htbp]
	\centering
	\begin{tikzpicture}
		\selectcolormodel{gray}
		\begin{axis}[
			axis lines = left, 
			xlabel = Размер, 
			ylabel = Время, 
			grid=both,
			minor grid style={gray!25}, 
			major grid style={gray!50},
			legend style={font=\scriptsize, legend pos=north west}
		]
			\addplot+[smooth, mark=none, thick] table [x index=0, y index=1] {../code/proctime.txt};
			\addlegendentry{Рекурсивный}
			\addplot+[smooth, mark=none, dashed] table [x index=0, y index=2] {../code/proctime.txt};
			\addlegendentry{Нерекурсивный}
		\end{axis}
	\end{tikzpicture}
	\caption{График зависимости процессорного времени от N (при k = 1000).}
\end{figure}

\section*{Вывод}
%<что сделали и получили в результате, кратко>

В данном разделе были проведены замеры процессорного времени для каждой реализации из предыдущего раздела. Полученные на основе этих замеров графики иллюстрируют рассчитанные ранее трудоёмкости рассматриваемых в работе алгоритмов.

Помимо этого, мной был проведён сравнительный анализ реализаций алгоритмов по критериям ёмкостной эффективности.

\clearpage
